% !TeX program = xelatex
% Build from this folder with:
%   latexmk -xelatex -interaction=nonstopmode -halt-on-error \
%     jc_landscape_and_parity_polytope.tex
\documentclass[11pt]{article}

\usepackage[margin=1in]{geometry}
\usepackage{fontspec}
\setmainfont{Times New Roman}
\setsansfont{Arial}
\usepackage{unicode-math}
\setmathfont{TeX Gyre Termes Math}
\usepackage{microtype}
\usepackage{mathtools}
\usepackage{amsthm}
\usepackage{booktabs}
\usepackage{array}
\usepackage{tabularx}
\usepackage{float}
\usepackage{needspace}
\usepackage{enumitem}
\usepackage{xcolor}
\usepackage{hyperref}
\usepackage{xurl}
\usepackage[capitalise,noabbrev]{cleveref}
\usepackage[
  backend=biber,
  style=ieee,
  sorting=none,
  maxbibnames=8
]{biblatex}
\addbibresource{references.bib}

\definecolor{linkblue}{HTML}{003D7C}
\hypersetup{
  colorlinks=true,
  linkcolor=linkblue,
  citecolor=linkblue,
  urlcolor=linkblue,
  pdftitle={Projected Gradient Descent for LDPC Decoding: From a Smooth Parity Penalty to Explicit Constraints},
  pdfauthor={BP--GD Mechanism Study}
}

\newtheorem{lemma}{Lemma}
\newtheorem{proposition}{Proposition}
\newtheorem{theorem}{Theorem}
\theoremstyle{remark}
\newtheorem{remark}{Remark}

\newcommand{\R}{\mathbb{R}}
\newcommand{\Cbox}{\mathcal{C}}
\newcommand{\Pc}{\mathcal{P}_c}
\newcommand{\Vc}{\mathcal{V}_c}
\newcommand{\proj}[2]{\Pi_{#1}\!\left(#2\right)}
\newcommand{\norm}[1]{\left\lVert#1\right\rVert}
\newcommand{\trans}{\mathsf{T}}

\title{Projected Gradient Descent for LDPC Decoding:\\
From a Smooth Parity Penalty to Explicit Constraints}
\author{BP--GD Mechanism Study}
\date{Technical research note --- 8 November 2026}

\begin{document}
\maketitle

\begin{abstract}
This paper explains a persistent performance gap between native exact
sum-product belief propagation (BP) and a family of projected-gradient
decoders for a length-$10{,}200$ $(3,6)$-regular low-density parity-check code
with $5{,}100$ parity checks.  On registered paired
cases, native projected gradient descent recovered $112/200$ transmitted
words at $6$~dB and $0/200$ at $1.5$~dB after 30 updates, whereas exact BP
recovered $200/200$ and $199/200$, respectively.  A later compound-box
projected-gradient construction reached $200/200$ after 90 updates on a
separate registered $6$-dB population, but it has no completed $1.5$-dB
transfer.  A related registered scalar control remained at $0/200$ at
$1.5$~dB after both 30 and 90 updates.

We show analytically how the smooth product parity term can assign zero cost
to both a correct local sign pattern and a pattern with two wrong signs, while a
continuous correction between them incurs positive cost.  A constructed
two-variable example then proves that an objective with the implemented term
shapes can possess a strict wrong local minimum on the box even when the
desired word has much lower objective value.  Exact projected stationarity
precludes escape by additional ordinary local descent, but it does not exclude
momentum, nonmonotone jumps, state-dependent matrices, different
initializations, different objectives, or expanded-variable methods.  We
therefore derive a differentiable convex parity cost based on distance to one
parity polytope per check and explain why finite quadratic penalties do not
generally impose exact feasibility.  On an exposed 200-case population at
$1.5$~dB, ordinary box-projected descent of one frozen instance of this cost
recovered $0/200$ words at every checkpoint through update 6000 and produced
no parity-valid readout; all 200 terminals remained above the registered
projected-gradient-mapping tolerance and are therefore optimization-unresolved,
not identified minimizers.  A replica variable for each check, together with
multipliers, can
enforce consistency without recreating the same product hill, although the
loopy intersection remains the fundamental-polytope relaxation and can have
fractional vertices.  The new count is an exposed-population mechanism
diagnostic, not a two-SNR confirmation or comparison with BP.
\end{abstract}

\section{Introduction}

Projected-gradient decoding is attractive because its operations are
differentiable, parallelizable, and compatible with parameter training.  Prior
work has formulated low-density parity-check (LDPC) decoding as nonconvex
projected-gradient optimization with smooth parity penalties and has used
restarts or learned parameters to address local minima
\cite{wadayama2019trainable}.  Linear-programming (LP) decoding takes another
route: it minimizes a linear channel cost over a relaxation derived from the
parity checks \cite{feldman2005lp}.  Efficient decomposition methods introduce
replica variables and Euclidean projection onto each parity polytope
\cite{barman2013decomposition}.

The present study began with a narrower engineering question: why does a
specific projected-gradient construction work well at $6$~dB yet remain far
behind exact sum-product BP at $1.5$~dB?  The campaign progressively separated
state, direction, step rule, check assignment, and update budget.  It found
that direct application of an exact recurrence can recover words where tested
true-gradient arms do not, but direct recurrence application is a fixed-point
update rather than descent of the declared objective.

The resulting claim is deliberately bounded.  We establish that the
implemented objective \emph{can} contain wrong local minima on the box and
that an exact projected zero cannot be escaped by further
ordinary projected descent.  We do not establish that every $1.5$-dB failure
has been registered as an exact local minimum, and we do not claim that every
method involving the partial derivatives of the objective function must fail.

The contributions are fourfold:

\begin{enumerate}[leftmargin=*,itemsep=2pt]
  \item We organize the controlled evidence without combining distinct case
        populations or treating a direction fit as a decoder result.
  \item We derive the partial derivatives at a hard vertex and isolate what
        happens when one check contains two wrong signs.
  \item We give a strict constructed local minimum on the box and state precisely
        what exact projected stationarity closes.
  \item We derive a differentiable convex parity cost, prove its exactness
        limitation, report an exposed-population diagnostic, and
        motivate one parity-polytope constraint per check with multipliers as
        a scientifically distinct successor formulation.
\end{enumerate}

\section{Controlled Evidence and Method Boundary}

All comparisons in the initial paired experiment held fixed the parity-check
matrix, transmitted word, noise realization, and vector of channel
log-likelihood ratios (LLRs).  A positive LLR favors bit zero.  Recovery means
that the hard-decision estimate equals the transmitted word; parity validity
alone is not counted as recovery.

Here \emph{native exact sum-product BP} means that the implementation evaluates
the declared sum-product recurrence without approximation.  It does not mean
exact maximum a posteriori or maximum-likelihood inference on the loopy Tanner
graph.

Exact sum-product BP retains directed variable-to-check and check-to-variable
messages and performs recipient exclusion.  The scalar projected-gradient
control retains one scalar $x_i\in[-1,1]$ for bit $i$, calculates the partial
derivatives of its objective function, applies its declared direction and step
rule, and projects componentwise onto the box.  These methods traverse the
same Tanner graph but do not retain or update the same quantities.

\begin{table}[H]
\centering
\caption{Registered recovery evidence.  Every entry is the number of
transmitted words recovered out of 200 cases.  Rows with different case
populations are not a single cross-SNR decoder evaluation.}
\label{tab:registered}
\small
\begin{tabularx}{\textwidth}{@{}>{\raggedright\arraybackslash}p{3.3cm}
  >{\raggedright\arraybackslash}p{2.6cm}
  >{\raggedright\arraybackslash}X
  >{\raggedright\arraybackslash}p{3.2cm}@{}}
\toprule
Experiment and population & SNR and budget & Projected-gradient or recurrence result & Exact BP or scope boundary \\
\midrule
Native composition, frames 60025--60224 & $6$~dB, 30 updates & Native projected GD: 112 & Exact BP: 200 \\
Native composition, frames 60225--60424 & $1.5$~dB, 30 updates & Native projected GD: 0 & Exact BP: 199 \\
Selected compound-box, frames 50025--50224 & $6$~dB, 30 / 90 updates & 177 / 200 & No completed exact $1.5$-dB transfer \\
Native compound budget diagnostic, frames 61745--61944 & $6$~dB, 30 / 90 updates & 113 / 193 & Separate population and construction \\
Native compound budget diagnostic, frames 61969--62168 & $1.5$~dB, 30 / 90 updates & 0 / 0 & Closes this frozen budget extension only \\
Convex parity-distance diagnostic, frames 61969--62168 & $1.5$~dB, 30 / 90 / 300 / 1000 / 3000 / 6000 updates & 0 / 0 / 0 / 0 / 0 / 0 & Exposed population; all 200 cap terminals unresolved; BP absent from candidate packet/source \\
State--direction--step factorial, frames 70921--71120 & $1.5$~dB, 30 updates & Free state + recurrence displacement + unit application: 196; every tested true-gradient cell: 0 & Successful cell is a fixed-point update, not a gradient step \\
\bottomrule
\end{tabularx}
\end{table}

The strongest retained gradient-based recovery result is the selected
compound-box construction at $6$~dB: $200/200$ transmitted words after 90
updates.  It is not a cross-SNR result approaching BP.  Conversely, the
$0/200$ registered result at $1.5$~dB through 90 updates closes the claim that
the frozen scalar control merely needed that larger cap; it is not a theorem
about every step sequence or every gradient-based decoder.

In the state--direction--step factorial, direct unit application of the
recurrence minus the current state recovered $196/200$ words at $1.5$~dB,
whereas all tested true-gradient cells recovered $0/200$.  This establishes an
update-operator boundary.  It does not establish that BP succeeds only by
increasing the objective, nor does endpoint nonincrease certify monotonicity
inside a finite update.

\section{The Implemented Smooth Objective}

Let $x\in[-1,1]^n$ denote the mean coordinate on which the compound objective
depends.  The implemented objective is

\begin{equation}
J_C(x)=
-\frac{L^{\trans}x}{2n}
+\frac{w_p}{n}\sum_c
  \left(1-\prod_{i\in c}x_i\right)^2
+\frac{w_b}{n}\sum_i(1-x_i^2)^2,
\label{eq:jc}
\end{equation}

where $w_p=2$ and $w_b=0.2$.  The production evaluator clips each LLR to
$[-20,20]$.  Define

\begin{equation}
p_c=\prod_{j\in c}x_j,
\qquad
q_{ci}=\prod_{\substack{j\in c\\j\ne i}}x_j.
\end{equation}

The exact partial derivative is

\begin{equation}
\partial_iJ_C(x)
=-\frac{L_i}{2n}
-\frac{2w_p}{n}\sum_{c\ni i}(1-p_c)q_{ci}
+\frac{4w_b}{n}x_i(x_i^2-1).
\label{eq:partial}
\end{equation}

The implementation may store directed-edge coordinates $z$, but the objective
depends only on their incident-edge means $x=Cz$.  Thus $J_C(z)=J_C(Cz)$, and
directions in $\ker C$ are exactly flat.  Claims about a well or barrier must
therefore be made in the mean coordinate, not by treating the full edge state
as an isolated minimum.

\subsection{Partial derivatives at a hard vertex}

\begin{lemma}[Partial derivatives at a hard vertex]
Let $s\in\{-1,+1\}^n$, and let $k_i$ be the number of violated checks incident
on bit $i$.  Then
\begin{equation}
\partial_iJ_C(s)
=-\frac{L_i}{2n}+\frac{4w_p}{n}k_i s_i.
\label{eq:vertex-force}
\end{equation}
\end{lemma}

\begin{proof}
At a satisfied check, $p_c=1$, so its contribution to
\cref{eq:partial} is zero.  At a violated check, $p_c=-1$ and
$q_{ci}=-s_i$, so its contribution is $4w_ps_i/n$.  The final term of
\cref{eq:partial} is zero at a hard vertex.  Summing the $k_i$ violated-check
contributions gives \cref{eq:vertex-force}.
\end{proof}

One violated check therefore applies the same parity magnitude to all six of
its incident coordinates at a hard vertex.  The complete partial derivative
can still distinguish those coordinates through the channel LLRs, their total
numbers $k_i$ of violated incident checks, overlapping checks, and, away from
the vertex, the partner products and the final term of \cref{eq:partial}.
Exact sum-product BP instead forms a check-to-variable message from the other
five incoming variable-to-check messages.  Recipient exclusion is the
structural contrast; \cref{eq:vertex-force} alone is not a decoder-performance
theorem.

The registered initialization is $x_i^{(0)}=\tanh(L_i/2)$ rather than a hard
vertex.  The lemma consequently explains a local mechanism, not the exact
initial trajectory.

\section{How Two Wrong Signs Create an Uphill Route}

At a hard-sign vector, a check containing zero wrong signs has product $+1$,
a check containing one wrong sign has product $-1$, and a check containing two
wrong signs again has product $+1$.  Hence one satisfied check assigns zero
parity penalty to both the all-correct pattern and a pattern with two wrong
signs.
This is not a defect in the parity test: both local patterns have even parity.
The channel evidence and overlapping checks must determine which local pattern
is globally preferable.

\begin{proposition}[Two-sign path through one check]
Hold all but two members of one check at $+1$ and set the remaining two to
$x_i=x_j=t$.  The contribution of that check is
\begin{equation}
J_{\mathrm{check}}(t)=\frac{w_p}{n}(1-t^2)^2.
\label{eq:pair-path}
\end{equation}
It is zero at $t=-1$ and $t=+1$, and it is positive for $|t|<1$.
\end{proposition}

\begin{proof}
Along the declared path, the check product is $t^2$.  Substitution into the
second term of \cref{eq:jc} gives \cref{eq:pair-path}.  In particular,
$J_{\mathrm{check}}(0)=w_p/n>0$.
\end{proof}

The contribution $(1-t^2)^2$ that favors hard magnitudes raises the same
intervening hill.  A gradient update can move every coordinate simultaneously;
the issue is not that it literally flips one bit at a time.  The issue is that
even a collective small correction can initially increase the objective.

This proposition exhibits one uphill path for one check term.  It does not
prove that every continuous path between two codewords crosses a positive rim
of the complete objective.  Such a claim requires minimizing the maximum
objective height over all connecting paths, together with a separately
declared analysis of finite jumps.

\subsection{A strict constructed local minimum on the box}

The following example proves possibility without claiming frequency on the
registered graph.  On $[-1,1]^2$, define

\begin{equation}
\begin{aligned}
F(x_1,x_2)
={}&-\frac{x_1+4.6x_2}{2}+2(1-x_1x_2)^2\\
&+0.2\left[(1-x_1^2)^2+(1-x_2^2)^2\right].
\end{aligned}
\label{eq:toy}
\end{equation}

The desired valid point has $F(+1,+1)=-2.8$.  Nevertheless, the following
holds.

\begin{proposition}[Constructed wrong local minimum on the box]
For the function in \cref{eq:toy},
\begin{equation}
x^*=(-1,-0.5),\qquad
F(x^*)=2.2625,\qquad
\nabla F(x^*)=(0.5,0).
\end{equation}
The point $x^*$ is a strict local minimum on $[-1,1]^2$.
\end{proposition}

\begin{proof}
At the active lower bound $x_1=-1$, any sufficiently small permitted change
has $d_1\ge0$.  Thus $\nabla F(x^*)^{\trans}d=0.5d_1>0$ whenever $d_1>0$.
When $d_1=0$, the remaining change is along $x_2$, where the second derivative
is
\begin{equation}
\frac{\partial^2 F}{\partial x_2^2}(x^*)=3.8>0.
\end{equation}
More strongly, the complete Hessian is
\begin{equation}
\nabla^2F(x^*)=
\begin{bmatrix}
2.6&0\\
0&3.8
\end{bmatrix}
\succ0.
\end{equation}
By continuity, the Hessian remains positive definite in a sufficiently small
neighborhood of $x^*$.  Taylor expansion then gives a nonnegative linear term
$0.5d_1$ for every permitted displacement and a strictly positive quadratic
term for every nonzero displacement.  Hence $x^*$ is a strict local minimum
on the box.
\end{proof}

For any $\eta>0$,

\begin{equation}
\proj{[-1,1]^2}{x^*-\eta\nabla F(x^*)}=x^*.
\end{equation}

The lower desired point exists, but an ordinary local projected-gradient step
cannot reach it from $x^*$.

\section{What Exact Projected Stationarity Proves}

Let $\Cbox=[-1,1]^n$.  For a differentiable function $J$, define

\begin{equation}
G_\eta(x)=
\frac{x-\proj{\Cbox}{x-\eta\nabla J(x)}}{\eta},
\qquad \eta>0.
\label{eq:gradient-residual}
\end{equation}

\begin{theorem}[Exact projected zero]
For $x^*\in\Cbox$ and $\eta>0$,
\begin{equation}
G_\eta(x^*)=0
\quad\Longleftrightarrow\quad
0\in\nabla J(x^*)+N_{\Cbox}(x^*),
\label{eq:kkt}
\end{equation}
where $N_{\Cbox}$ is the convex normal cone.  Consequently,
$\nabla J(x^*)^{\trans}d\ge0$ for every $d$ such that
$x^*+\epsilon d\in\Cbox$ for all sufficiently small $\epsilon>0$, and
\begin{equation}
\proj{\Cbox}{x^*-\alpha\nabla J(x^*)}=x^*
\qquad\text{for every }\alpha>0.
\end{equation}
\end{theorem}

\begin{proof}
The Euclidean projection characterization gives
$x^*=\proj{\Cbox}{x^*-\eta\nabla J(x^*)}$ if and only if
$-\eta\nabla J(x^*)\in N_{\Cbox}(x^*)$.  Because the normal cone is a cone,
this is equivalent to \cref{eq:kkt} and is unchanged by any positive scalar
$\alpha$.  The final inequality follows directly from the definition of the
normal cone.
\end{proof}

The frozen registered scalar protocol uses

\begin{equation}
D=\frac{13.6}{n}I,
\qquad
G_\eta^D(x)=
\frac{x-\proj{\Cbox}{x-\eta D^{-1}\nabla J_C(x)}}{\eta}.
\end{equation}

Since $D$ is a positive scalar multiple of the identity, this is the same
Euclidean direction with a rescaled step.  The exact zero sets coincide after
the corresponding scalar conversion.  Finite numerical tolerances do not:
for $n=10200$, $D^{-1}$ scales an interior partial derivative by $750$.

An exact projected zero therefore closes additional ordinary local
$-\nabla J_C$ descent from that state.  It does not close a route that avoids
the basin, carries nonzero auxiliary velocity, takes a finite nonmonotone
jump, changes the initialization, applies a nonscalar or state-dependent
matrix before stationarity, changes the objective, or introduces new
variables.

\subsection{Registered and indicative evidence}

The registered result is $0/200$ recovery at $1.5$~dB through 90 updates for
the frozen native compound control.  It does not certify exact stationarity.

An external NumPy continuation on the same 200 exposed $1.5$-dB frames
reported $0/200$ recoveries, 197 terminal states meeting its own unscaled
threshold $\|G_\eta(x)\|_\infty\le 10^{-8}$, three 6000-update caps, a median
of $744.5$ wrong bits, and
$J_C(x_{\mathrm{word}})<J_C(x_{\mathrm{terminal}})$ on all 200 frames.  This
is useful indicative mechanism evidence, not a registered result.  Its finite
tolerance differs in scale from the registered $D^{-1}$ diagnostic, and it did
not perform the registered independent $10^{-10}$ confirmation pass.

The defensible interpretation is therefore that the external continuation is
consistent with trajectories ending in wrong stationary regions.  It does not
certify local minimality.  The registered result alone establishes only that
the frozen method did not recover any of those 200 transmitted words within
90 updates.

\section{A Convexity Boundary for Exact Parity Penalties}

Parity can be placed in a differentiable cost.  The obstruction is to asking
one ordinary penalty with a finite coefficient to be simultaneously smooth,
convex, exact,
zero only at binary parity-valid points, and easy to minimize by ordinary
gradient descent.

Let $\Vc$ be the set of binary even-parity patterns for check $c$, and let
$\Pc=\operatorname{conv}(\Vc)$ be its parity polytope
\cite{barman2013decomposition}.

\begin{theorem}[Convex zero-set obstruction]
Suppose $p$ is nonnegative and convex and satisfies $p(v)=0$ for every
$v\in\Vc$.  Then $p(u)=0$ for every $u\in\Pc$.
\end{theorem}

\begin{proof}
Every $u\in\Pc$ has a representation
$u=\sum_{v\in\Vc}\alpha_vv$, where $\alpha_v\ge0$ and
$\sum_v\alpha_v=1$.  Convexity and nonnegativity give
\begin{equation}
0\le p(u)
\le\sum_{v\in\Vc}\alpha_vp(v)=0.
\end{equation}
Hence $p(u)=0$ throughout $\Pc$.
\end{proof}

A convex penalty cannot therefore have precisely the multiple isolated binary
even-parity patterns as its zero-valued global minima.  The product
construction obtains isolated hard zeros by giving up convexity, thereby
permitting multiple local minima and intervening hills.

\section{A Differentiable Convex Parity Cost}

Let $b\in[0,1]^n$ be a shared vector, let $S_cb$ extract the entries incident
on check $c$, and let $f_{\mathrm{sig}}$ be a differentiable signal cost.
Consider

\begin{equation}
J_\rho(b)=
f_{\mathrm{sig}}(b)
+\frac{\rho}{2}\sum_c
\norm{S_cb-\proj{\Pc}{S_cb}}^2.
\label{eq:distance-cost}
\end{equation}

\begin{proposition}[Squared distance to the parity polytope]
If every $\Pc$ is nonempty, closed, and convex, then the parity part of
\cref{eq:distance-cost} is convex and continuously differentiable.  Moreover,
\begin{equation}
\nabla J_\rho(b)=
\nabla f_{\mathrm{sig}}(b)
+\rho\sum_c S_c^{\trans}
\left[S_cb-\proj{\Pc}{S_cb}\right].
\label{eq:distance-gradient}
\end{equation}
If $f_{\mathrm{sig}}$ is convex, then $J_\rho$ is convex.
\end{proposition}

\begin{proof}
For a nonempty closed convex set $P$, the function
$\tfrac12\operatorname{dist}(u,P)^2$ is convex and continuously
differentiable, with gradient $u-\Pi_P(u)$.  Applying the chain rule to each
linear extraction $S_c$ yields \cref{eq:distance-gradient}.  Convexity is
preserved by linear composition, nonnegative summation, and addition of
$f_{\mathrm{sig}}$.
\end{proof}

This parity contribution pulls each local configuration toward its parity
polytope instead of creating a product hill between valid hard
configurations.  It does not guarantee a binary output: it is zero throughout
each $\Pc$, including at fractional mixtures.  Nor is the finite quadratic
penalty generally exact.

The scalar example

\begin{equation}
\min_x x\quad\text{subject to}\quad x=0
\end{equation}

has constrained solution $x=0$.  Replacing the constraint by

\begin{equation}
F_\rho(x)=x+\frac{\rho}{2}x^2
\end{equation}

gives $x_\rho=-1/\rho$, which violates the constraint for every finite
$\rho$.  Increasing $\rho$ reduces the violation without making it exactly
zero.  In \cref{eq:distance-cost}, the restoring term in
\cref{eq:distance-gradient} vanishes at feasibility, so a nonzero signal
gradient can balance it at a slightly infeasible point.

\paragraph{Exposed-population diagnostic.}
Experiment \path{RQ2-CONVEX-PARITY-DISTANCE-001} evaluated one frozen
instance of \cref{eq:distance-cost} on the already exposed $1.5$-dB frames
61969--62168.  It retained one scalar $b_i\in[0,1]$ per bit, used the linear
signal term with $\rho=96$, calculated the partial derivatives through exact
Euclidean projection onto every degree-six parity polytope, applied the fixed
step $1/288$, and projected componentwise onto the box.  Its frozen source
defines no sigmoid, restart, momentum, multiplier, or state-dependent matrix,
and its candidate packet contains no BP quantity.

The latched transmitted-word recovery count was $0/200$ at updates 30, 90,
300, 1000, 3000, and 6000.  No parity-valid hard-decision estimate occurred at
any update.  No stored candidate state, including every checkpoint and
terminal, had a readout tie.  Every case reached the update-6000 cap with
$\norm{G(b)}_\infty>10^{-8}$, so all 200 terminals are
\texttt{optimization\_unresolved}.  The terminal mapping-norm median was
$2.2772885202382298\mathbin{\times}10^{-4}$, with range
$2.4881\mathbin{\times}10^{-5}$ to $0.0115097$.  After the post-commit scorer
opened the transmitted word, the terminal hard-decision estimates had median
382 coded-bit errors (range 148--595) and median syndrome weight 387 (range
224--508).  The soft states had median 3274 fractional coordinates and median
maximum local check distance $0.102299$.  Primary and replay scientific
traces were identical.

Although the objective values were essentially flat late in the run, none of
the cases met the frozen numerical-stationarity tolerance.  Convexity
therefore does not license a claim that these terminals are global minimizers,
and the result does not close this objective.  The source-only follow-up below
tested faster same-objective convergence without clearing its population-run
gate.  The exact constrained replica-and-multiplier formulation and the
untouched two-hop branch were not tested.

\paragraph{Source-only solver-successor gates.}
Production-size synthetic gates on the same frozen objective did not reach
$\norm{G_{288}}_\infty\leq 10^{-8}$, so no second $1.5$-dB population run was
launched.  The feasible accelerated arm ended after 6000 updates at
$\norm{G_{288}}_\infty=1.8247366\mathbin{\times}10^{-3}$; its synthetic suite
passed 44,031/44,031 assertions.  A hybrid of those 6000 updates and 30,000
monotone local-curvature updates ended at
$\norm{G_{288}}_\infty=2.3504324\mathbin{\times}10^{-4}$ after 73,511
evaluations and 43,510 backtracks.  Optim.jl 2.3.2 \texttt{LBFGSB}, capped at
5000 iterations, ended at
$\norm{G_{288}}_\infty=5.4017116\mathbin{\times}10^{-5}$ for $m=10$ and
$6.3955257\mathbin{\times}10^{-5}$ for $m=50$.  The local-curvature suite
passed 15,233/15,233 assertions.

These are synthetic source-feasibility diagnostics, not $1.5$-dB population
evidence and not evidence that the objective fails.  A decisive
global-minimum audit now requires a materially stronger certified optimizer,
or a relaxed tolerance paired with an explicit error bound.  The active
two-hop branch is untouched.

\section{Replica Variables and Multipliers}

Exact consistency can instead remain a constraint.  Introduce one replica
variable $z_c$ for each check and impose

\begin{equation}
z_c=S_cb,
\qquad
z_c\in\Pc.
\label{eq:replicas}
\end{equation}

With residual $r_c=S_cb-z_c$ and multiplier $y_c$, the augmented Lagrangian is

\begin{equation}
\mathcal{L}_\rho(b,z,y)
=f_{\mathrm{sig}}(b)
+\sum_c\left[
y_c^{\trans}r_c+\frac{\rho}{2}\norm{r_c}^2
\right],
\qquad z_c\in\Pc.
\label{eq:augmented}
\end{equation}

Its partial derivative with respect to the shared vector is

\begin{equation}
\nabla_b\mathcal{L}_\rho
=\nabla f_{\mathrm{sig}}(b)
+\sum_cS_c^{\trans}(y_c+\rho r_c).
\label{eq:augmented-gradient}
\end{equation}

At exact consistency, $r_c=0$, the quadratic contribution disappears but the
multiplier need not.  The first-order balance can remain

\begin{equation}
\nabla f_{\mathrm{sig}}(b)
+\sum_cS_c^{\trans}y_c=0.
\end{equation}

Thus the multipliers can retain a nonzero normal force while the shared and
replica variables agree.  A typical alternating direction method updates
the shared variables, projects each replica variable onto its parity
polytope, and updates the multipliers from the residuals
\cite{boyd2011admm,barman2013decomposition}.  This uses partial derivatives,
but it is not ordinary descent of the unconstrained scalar objective $J_C$.
It changes the variables, feasible sets, and update operation.

\section{Exactness Boundary on a Loopy Tanner Graph}

Each $\Pc$ exactly describes the convex hull of the valid binary patterns for
one parity check.  On a loopy Tanner graph, however, intersecting these local
polytopes through the shared variables gives the fundamental-polytope
relaxation used in LP decoding \cite{feldman2005lp}.  That relaxation can have
fractional vertices, commonly called pseudocodewords in the LP-decoding
literature.

Consequently,

\begin{equation}
\begin{gathered}
\text{exact feasibility in every check's parity polytope}\\
\text{does not guarantee a binary codeword}.
\end{gathered}
\end{equation}

A linear cost over the exact global codeword polytope has an optimal codeword
vertex.  A unique optimum selects that vertex.  If several codewords tie,
fractional mixtures on the optimal face are also optimal and require a
deterministic vertex or tie rule.  Describing or projecting onto the exact
global polytope can be substantially more expensive than using the relaxation
formed from the individual checks.  The difficulty therefore moves from a
nonconvex product landscape
into the representation and enforcement of the feasible set; it does not
disappear.

\Needspace{8\baselineskip}
\section{Prospective Research Direction}

The mathematically distinct successor formulation is

\begin{equation}
\boxed{
\begin{aligned}
&\text{linear or differentiable lifted signal cost}\\
&\quad+\text{ parity-polytope feasibility}\\
&\quad+\text{ shared-to-local consistency and multipliers}
\end{aligned}}
\end{equation}

The finite differentiable penalty in \cref{eq:distance-cost} now has the
exposed-population diagnostic reported above, but that run stopped above its
mapping tolerance and does not identify its global minimizers.  The expanded
constrained formulation in this section still has no recovery evidence in the
present repository and is not introduced as a validated decoder.  A
prospective evaluation would need to declare its state dimension, projection
algorithm, multiplier update, initialization, stopping and readout rules,
deterministic replay checks, operation and memory budgets, and recovery and
coded-bit-error metrics before opening held-out cases.

It also does not supersede the separate state-dependent two-hop matrix branch.
That branch applies a bounded matrix to the partial derivatives of $J_C$ and
has a retained teacher-labelled fit, but no fresh direction-screen or decoder
result.  It must be completed or amended under its own registered gates.

\section{Conclusion}

The observed decoding failure need not mean that the calculated partial
derivatives are wrong.  The mathematics establishes a more specific
possibility: a descent method can faithfully optimize an objective term
structure that permits a wrong local minimum on the box.  A satisfied check
cannot distinguish the
all-correct hard pattern from a hard pattern with two wrong signs, and the smooth product term can
rise along a collective correction between them.  The constructed example
proves that this term structure can yield a strict wrong local minimum on the
box even when the desired word has lower cost.

The experimental conclusion remains narrower.  The frozen scalar protocol
recovered $0/200$ words at $1.5$~dB through 90 registered updates.  The
external continuation is consistent with trajectories ending in wrong
stationary regions, but it is not a registered exact-stationarity certificate.
The evidence closes the frozen scalar step grid and normalization against the
specific budget extension from 30 to 90 updates.  As a research-priority
decision, further repetitions that preserve the same state, objective,
direction, and projection while changing only scalar implementation details
should be deprioritized.  Arbitrary discrete scalar-step schedules remain
mathematically unclosed, and the evidence does not justify abandoning every
use of partial derivatives.

The exposed convex parity-distance diagnostic removed nonglobal local wells
from the tested objective but recovered $0/200$ words through update 6000 and
produced no parity-valid hard-decision estimate.  This is not yet evidence
that the convex objective has the wrong global minimizer: every case stopped
at the cap above the frozen mapping tolerance.  The source-only successor
gates also remained above tolerance, so no second population run was launched.
The immediate question is therefore convergence, not another interpretation
of the unresolved terminals.  A decisive global-minimum audit requires a
materially stronger certified optimizer, or a relaxed tolerance paired with
an explicit error bound; only then can the signal cost, finite penalty, and
fundamental-polytope relaxation be classified.

The constrained formulation changes what differentiation is asked to do.  A
differentiable signal cost ranks candidate explanations; parity polytopes,
projections, and multipliers enforce the coding structure.  Whether that
expanded constrained formulation can approach BP at both $6$~dB and
$1.5$~dB remains a prospective experimental question.

\section*{Reproducibility and Evidence Boundary}

This paper is derived from the repository's explanatory and analytical
records.  The implemented objective is in
\path{../implementation/two_hop_metric_action/compound_objective_jc.jl}.
The full projected-descent derivation and result classification are in
\path{../research/jc_projected_gradient_boundary.md}.  The registered but
unfinished continuation card is
\path{../research/rq2_jc_projected_stationarity_001_card.md}.  The active
two-hop branch is scoped in \path{../research/active_goal.md}.  The convex
parity-distance diagnostic is specified in
\path{../research/rq2_convex_parity_distance_001_card.md}, implemented in
\path{../implementation/convex_parity_distance/}, and committed under
\path{../results/rq2_convex_parity_distance_001/2026-11-08_v1/}.  Its
candidate and replay scientific-trace SHA-256 values both equal
\path{2020306bb9d09d20e025528c2b9eba8a398f67e8c896b7329c8b3411c1c1dd20}.
It used only the already exposed frames 61969--62168; the candidate packet
contained no BP quantity, the frozen candidate source contained no BP call,
and the run did not open the untouched confirmation ranges.

\clearpage
\printbibliography

\end{document}
